%! Vector-Valued Functions — Unit 4, Lesson 4.9 — Exit Ticket (Answer Key)
\documentclass[10pt]{article}
\usepackage{precalculus-article}
\usepackage{precalculus-key}   % pulls in -boxes; do NOT also load -boxes
\newcommand{\TallMath}[1]{$\displaystyle #1\rule[-1.4em]{0pt}{3.2em}$}

\begin{document}

\pageheader{Unit 4, Lesson 4.9}{Exit Ticket --- Answer Key}
\namedateperiod

\vspace{0.12in}

\begin{notesbox}{Vector-Valued Functions Check --- Key}

A particle moves with position vector $\vec{p}(t)=\langle t^{2}-1,\;2t+3\rangle$
(distances in meters, $t$ in seconds).

\vspace{10pt}
\textbf{1.}\quad Find $\vec{p}(2)$.  What does this vector represent in context?

\vspace{6pt}
$\vec{p}(2) = $ \ans{$\langle 3,\;7\rangle$}

\vspace{6pt}
In context: \ans{The particle is at position $(3,\,7)$ meters from the origin
at $t=2$ seconds.}

\begin{teachernote}
  Students sometimes confuse the vector $\langle 3,7\rangle$ with the scalar
  distance. Reinforce: the \emph{vector} gives the coordinates of position;
  the \emph{magnitude} (Question 2) gives the scalar distance from the origin.
\end{teachernote}

\vspace{12pt}
\textbf{2.}\quad $\|\vec{p}(2)\| = \sqrt{3^{2}+7^{2}} = \sqrt{9+49} = \sqrt{58} \approx $ \ans{$7.6$ m}

\vspace{12pt}
\textbf{3.}\quad $\vec{v}(1) = \langle 2(1),\;2\rangle = $ \ans{$\langle 2,\;2\rangle$} \qquad
Speed $= \sqrt{4+4} = \sqrt{8} = $ \ans{$2\sqrt{2} \approx 2.8$ m/s}

\end{notesbox}

\end{document}
